\documentclass[13pt,a4paper]{extarticle}
\usepackage{fontspec}
\usepackage{unicode-math}
\setmainfont{Latin Modern Roman}
\setmathfont{Latin Modern Math}
\usepackage{amsmath}
\usepackage[margin=1.8cm]{geometry}
\usepackage{enumitem}
\usepackage{array}
\usepackage{multicol}
\setlength{\columnsep}{1.2em}
\usepackage{eurosym}
\linespread{1.04}
\renewcommand{\arraystretch}{1.28}
\usepackage{newunicodechar}
\newunicodechar{·}{\kern-.13em\textperiodcentered\kern-.13em}
\usepackage{parskip}
\setlength{\parindent}{0pt}
\begin{document}
\section*{Exercicis i problemes}

\begin{quote}\small Llistat d'exercicis per practicar i consolidar els continguts del bloc. Segueix l'ordre dels apunts i inclou exercicis de diferents nivells de dificultat.\end{quote}


\vspace{1mm}\hrule\vspace{3mm}


\section*{1. Polinomis}

\par\noindent\begin{minipage}{\linewidth}
\textbf{1.} Per a cada polinomi, indica'n el grau, el coeficient principal i el terme independent:\par

\begin{tabular}{@{}p{.46\linewidth}p{.46\linewidth}@{}}
\textbf{a)} $3x^4-2x^3+x-7$ & \textbf{b)} $-5x^3+4x^2-1$ \\[2pt]
\textbf{c)} $x^5-3x^2$ & \textbf{d)} $8$ \\[2pt]
\end{tabular}\par
\end{minipage}\par\vspace{1.6mm}
\par\noindent\begin{minipage}{\linewidth}
\textbf{2.} Avalua cada polinomi en el valor indicat i digues si es tracta d'una arrel:\par

\begin{itemize}[leftmargin=0pt,itemsep=3pt,topsep=3pt,parsep=0pt,label={}]
\item \textbf{a)} $p(x)=x^3-4x+3$ en $x=1$
\item \textbf{b)} $p(x)=x^3-4x+3$ en $x=2$
\item \textbf{c)} $q(x)=2x^2-5x-3$ en $x=3$
\item \textbf{d)} $q(x)=2x^2-5x-3$ en $x=-1$
\end{itemize}
\end{minipage}\par\vspace{1.6mm}
\par\noindent\begin{minipage}{\linewidth}
\textbf{3.} Donats $p(x)=2x^3-x^2+4x-1$ i $q(x)=-x^3+3x^2-2x+5$, calcula $p(x)+q(x)$ i $p(x)-q(x)$.\par
\end{minipage}\par\vspace{1.6mm}
\par\noindent\begin{minipage}{\linewidth}
\textbf{4.} Desenvolupa:\par

\begin{tabular}{@{}p{.46\linewidth}p{.46\linewidth}@{}}
\textbf{a)} $(2x-1)(x^2+3x-2)$ & \textbf{b)} $(x+3)^2$ \\[2pt]
\textbf{c)} $(x-2)(x+2)$ &  \\[2pt]
\end{tabular}\par
\end{minipage}\par\vspace{1.6mm}
\par\noindent\begin{minipage}{\linewidth}
\textbf{5.} Divideix i indica el quocient i el residu:\par

\begin{itemize}[leftmargin=0pt,itemsep=3pt,topsep=3pt,parsep=0pt,label={}]
\item \textbf{a)} $(2x^3-3x^2+4x-5):(x^2-1)$
\item \textbf{b)} $(x^4-1):(x^2+1)$
\end{itemize}
\end{minipage}\par\vspace{1.6mm}
\par\noindent\begin{minipage}{\linewidth}
\textbf{6.} Divideix per Ruffini, indicant el quocient i el residu:\par

\begin{itemize}[leftmargin=0pt,itemsep=3pt,topsep=3pt,parsep=0pt,label={}]
\item \textbf{a)} $(x^3-6x^2+11x-6):(x-2)$
\item \textbf{b)} $(2x^3+x^2-5x+2):(x+2)$
\item \textbf{c)} $(x^4-5x^2+4):(x-1)$
\end{itemize}
\end{minipage}\par\vspace{1.6mm}
\vspace{1mm}\hrule\vspace{3mm}

\section*{2. Factorització de polinomis}

\par\noindent\begin{minipage}{\linewidth}
\textbf{7.} Factoritza completament:\par

\begin{tabular}{@{}p{.46\linewidth}p{.46\linewidth}@{}}
\textbf{a)} $x^3-5x^2+6x$ & \textbf{b)} $2x^3+2x^2-12x$ \\[2pt]
\textbf{c)} $x^4-9x^2$ & \textbf{d)} $3x^3-12x$ \\[2pt]
\end{tabular}\par
\end{minipage}\par\vspace{1.6mm}
\par\noindent\begin{minipage}{\linewidth}
\textbf{8.} Factoritza completament (totes les arrels són enteres):\par

\begin{itemize}[leftmargin=0pt,itemsep=3pt,topsep=3pt,parsep=0pt,label={}]
\item \textbf{a)} $x^3-2x^2-5x+6$
\item \textbf{b)} $x^3+2x^2-5x-6$
\end{itemize}
\end{minipage}\par\vspace{1.6mm}
\par\noindent\begin{minipage}{\linewidth}
\textbf{9.} Factoritza completament (queda un factor de grau $2$ irreductible):\par

\begin{itemize}[leftmargin=0pt,itemsep=3pt,topsep=3pt,parsep=0pt,label={}]
\item \textbf{a)} $x^3-3x^2+4x-2$
\item \textbf{b)} $x^3+x^2+x-3$
\end{itemize}
\end{minipage}\par\vspace{1.6mm}
\par\noindent\begin{minipage}{\linewidth}
\textbf{10.} Factoritza completament (les arrels del factor de grau $2$ són reals però irracionals):\par

\begin{itemize}[leftmargin=0pt,itemsep=3pt,topsep=3pt,parsep=0pt,label={}]
\item \textbf{a)} $x^3-2x^2-2x+4$
\item \textbf{b)} $x^3-x^2-3x+3$
\end{itemize}
\end{minipage}\par\vspace{1.6mm}
\par\noindent\begin{minipage}{\linewidth}
\textbf{11.} Factoritza completament (el polinomi no és mònic: treu primer el factor comú numèric):\par

\begin{itemize}[leftmargin=0pt,itemsep=3pt,topsep=3pt,parsep=0pt,label={}]
\item \textbf{a)} $2x^3-4x^2-32x+64$
\item \textbf{b)} $3x^3+3x^2-12x-12$
\end{itemize}
\end{minipage}\par\vspace{1.6mm}
\par\noindent\begin{minipage}{\linewidth}
\textbf{12.} Factoritza completament cada polinomi (no hi ha cap indicació de quin cas és cadascun):\par

\begin{itemize}[leftmargin=0pt,itemsep=3pt,topsep=3pt,parsep=0pt,label={}]
\item \textbf{a)} $x^3-7x+6$
\item \textbf{b)} $x^3-4x^2+6x-4$
\item \textbf{c)} $x^3+x^2-3x-3$
\item \textbf{d)} $2x^3-3x^2-11x+6$
\item \textbf{e)} $3x^3-5x^2-4x+4$
\end{itemize}
\end{minipage}\par\vspace{1.6mm}
\par\noindent\begin{minipage}{\linewidth}
\textbf{13.} Factoritza completament $x^5-3x^4-5x^3+15x^2+4x-12$ (totes les arrels són enteres).\par
\end{minipage}\par\vspace{1.6mm}
\par\noindent\begin{minipage}{\linewidth}
\textbf{14.} Factoritza completament $x^6+2x^5-5x^4-7x^3-2x^2+5x+6$ (quatre arrels són enteres i queda un factor de grau $2$ irreductible).\par
\end{minipage}\par\vspace{1.6mm}
\par\noindent\begin{minipage}{\linewidth}
\textbf{15.} \emph{(guiat)} Comprova, multiplicant, que $x^4+4=(x^2-2x+2)(x^2+2x+2)$, i comprova després que els dos factors de grau $2$ són irreductibles.\par
\end{minipage}\par\vspace{1.6mm}
\par\noindent\begin{minipage}{\linewidth}
\textbf{16.} Calcula el MCD i el MCM de cada parell de polinomis:\par

\begin{itemize}[leftmargin=0pt,itemsep=3pt,topsep=3pt,parsep=0pt,label={}]
\item \textbf{a)} $A(x)=x^2-9$, $B(x)=x^2-x-6$
\item \textbf{b)} $A(x)=x^3-x$, $B(x)=x^2+2x+1$
\end{itemize}
\end{minipage}\par\vspace{1.6mm}
\par\noindent\begin{minipage}{\linewidth}
\textbf{17.} Troba un polinomi que tingui exactament les arrels indicades:\par

\begin{itemize}[leftmargin=0pt,itemsep=3pt,topsep=3pt,parsep=0pt,label={}]
\item \textbf{a)} $2$, $-1$ i $3$
\item \textbf{b)} $0$ i $-2$ (doble)
\end{itemize}
\end{minipage}\par\vspace{1.6mm}
\vspace{1mm}\hrule\vspace{3mm}

\section*{3. Fraccions algebraiques}

\par\noindent\begin{minipage}{\linewidth}
\textbf{18.} Simplifica:\par

\begin{tabular}{@{}p{.46\linewidth}p{.46\linewidth}@{}}
\textbf{a)} $\dfrac{x^2-1}{x^2+2x+1}$ & \textbf{b)} $\dfrac{x^2-4}{x^2-5x+6}$ \\[2pt]
\textbf{c)} $\dfrac{x^2+3x}{x^2-9}$ & \textbf{d)} $\dfrac{x^3-x}{x^2-x}$ \\[2pt]
\end{tabular}\par
\end{minipage}\par\vspace{1.6mm}
\par\noindent\begin{minipage}{\linewidth}
\textbf{19.} Digues si cada igualtat és certa o falsa:\par

\begin{tabular}{@{}p{.46\linewidth}p{.46\linewidth}@{}}
\textbf{a)} $\dfrac{x+2}{x^2+4}=\dfrac{1}{x+2}$ & \textbf{b)} $\dfrac{x-3}{x^2-9}=\dfrac{1}{x+3}$ \\[2pt]
\textbf{c)} $\dfrac{x^2+2x}{2x}=x+2$ & \textbf{d)} $\dfrac{x^2-2x}{x-2}=x$ \\[2pt]
\end{tabular}\par
\end{minipage}\par\vspace{1.6mm}
\par\noindent\begin{minipage}{\linewidth}
\textbf{20.} Opera i simplifica:\par

\begin{itemize}[leftmargin=0pt,itemsep=3pt,topsep=3pt,parsep=0pt,label={}]
\item \textbf{a)} $\dfrac{1}{x-1}+\dfrac{2}{x+1}$
\item \textbf{b)} $\dfrac{x}{x^2-4}-\dfrac{1}{x+2}$
\item \textbf{c)} $\dfrac{3}{x}+\dfrac{x}{x-3}$
\end{itemize}
\end{minipage}\par\vspace{1.6mm}
\par\noindent\begin{minipage}{\linewidth}
\textbf{21.} Opera i simplifica:\par

\begin{itemize}[leftmargin=0pt,itemsep=3pt,topsep=3pt,parsep=0pt,label={}]
\item \textbf{a)} $\dfrac{x+3}{x-1}\cdot\dfrac{x^2-1}{(x+3)^2}$
\item \textbf{b)} $\dfrac{x^2-4}{x+1}:\dfrac{x-2}{x+1}$
\end{itemize}
\end{minipage}\par\vspace{1.6mm}
\par\noindent\begin{minipage}{\linewidth}
\textbf{22.} Opera i simplifica: $\dfrac{1}{x}+\dfrac{1}{x+1}-\dfrac{2}{x^2+x}$.\par
\end{minipage}\par\vspace{1.6mm}
\par\noindent\begin{minipage}{\linewidth}
\textbf{23.} Demostra les identitats següents:\par

\begin{itemize}[leftmargin=0pt,itemsep=3pt,topsep=3pt,parsep=0pt,label={}]
\item \textbf{a)} $\dfrac{1}{x-1}-\dfrac{1}{x+1}=\dfrac{2}{x^2-1}$
\item \textbf{b)} $\dfrac{x}{x+1}+\dfrac{1}{x+1}=1$
\end{itemize}
\end{minipage}\par\vspace{1.6mm}
\vspace{1mm}\hrule\vspace{3mm}

\section*{4. Problemes en context}

\par\noindent\begin{minipage}{\linewidth}
\textbf{24.} \emph{(aplicat)} Un rectangle té una àrea de $x^3-x^2-4x+4$ i un dels costats fa $x-1$. Troba, factoritzant, l'expressió de l'altre costat.\par
\end{minipage}\par\vspace{1.6mm}
\par\noindent\begin{minipage}{\linewidth}
\textbf{25.} \emph{(aplicat)} Un dipòsit prismàtic té un volum $V(x)=x^3+2x^2-5x-6$. Si dues de les seves dimensions són $x+1$ i $x-2$, troba la tercera.\par
\end{minipage}\par\vspace{1.6mm}
\par\noindent\begin{minipage}{\linewidth}
\textbf{26.} \emph{(aplicat)} Simplifica $\dfrac{x^2-9}{x-3}$ i avalua el resultat per a $x=5$.\par
\end{minipage}\par\vspace{1.6mm}
\par\noindent\begin{minipage}{\linewidth}
\textbf{27.} Factoritza i simplifica $\dfrac{x^3-4x}{x^2-x-2}$.\par
\end{minipage}\par\vspace{1.6mm}
\par\noindent\begin{minipage}{\linewidth}
\textbf{28.} Factoritza, simplifica i indica els valors de $x$ pels quals l'expressió original no està definida: $\dfrac{x^3-2x^2-3x}{x^2-9}$.\par
\end{minipage}\par\vspace{1.6mm}
\par\noindent\begin{minipage}{\linewidth}
\textbf{29.} \emph{(aplicat)} D'una làmina de cartró rectangular de $20\,\text{cm}\times12\,\text{cm}$ es retalla un quadrat de costat $x$ a cada cantonada i es doblega per formar una caixa sense tapa.\par

\begin{itemize}[leftmargin=0pt,itemsep=3pt,topsep=3pt,parsep=0pt,label={}]
\item \textbf{a)} Expressa el volum $V(x)$ de la caixa com un polinomi.
\item \textbf{b)} Calcula el volum per a $x=2\,\text{cm}$.
\end{itemize}
\end{minipage}\par\vspace{1.6mm}
\par\noindent\begin{minipage}{\linewidth}
\textbf{30.} \emph{(aplicat)} Una empresa estima que, en fabricar $x$ centenars d'unitats d'un producte, els ingressos (en milers d'€) són $I(x)=x^3-x^2-2x+8$ i els costos són $C(x)=x^2+3x+2$.\par

\begin{itemize}[leftmargin=0pt,itemsep=3pt,topsep=3pt,parsep=0pt,label={}]
\item \textbf{a)} Calcula el benefici $B(x)=I(x)-C(x)$.
\item \textbf{b)} Factoritza $B(x)$ i troba els punts d'equilibri (valors de $x>0$ pels quals el benefici és nul).
\item \textbf{c)} Avalua $B(2)$ i interpreta el signe del resultat.
\end{itemize}
\end{minipage}\par\vspace{1.6mm}
\par\noindent\begin{minipage}{\linewidth}
\textbf{31.} El nombre d'usuaris (en milers) d'una aplicació mòbil, $t$ mesos després del llançament, es modelitza per $N(t)=t^3-5t^2+2t+8$.\par

\begin{itemize}[leftmargin=0pt,itemsep=3pt,topsep=3pt,parsep=0pt,label={}]
\item \textbf{a)} Calcula $N(5)$.
\item \textbf{b)} Factoritza $N(t)$ i indica per a quins valors de $t\geq0$ el model prediu zero usuaris.
\end{itemize}
\end{minipage}\par\vspace{1.6mm}
\vspace{1mm}\hrule\vspace{3mm}

\section*{5. Exercicis inspirats en les PAU}

\par\noindent\begin{minipage}{\linewidth}
\emph{\small Aquests exercicis estan inspirats en proves PAU (Selectivitat) dels darrers anys (2022–2026). És important tenir en compte que, als exàmens originals, aquests càlculs formen part d'un problema més ampli —normalment de càlcul, matrius o geometria, propis de 2n de Batxillerat— i aquí n'aïllem només el fragment o procediment algebraic (factorització, equacions, fraccions algebraiques...) que té relació directa amb aquest tema.}\par
\end{minipage}\par\vspace{1.6mm}
\par\noindent\begin{minipage}{\linewidth}
\textbf{32.} \emph{(PAU 2024)} Factoritza completament $x^3-4x^2+4x$ i troba'n totes les arrels.\par
\end{minipage}\par\vspace{1.6mm}
\par\noindent\begin{minipage}{\linewidth}
\textbf{33.} \emph{(PAU 2026)} Factoritza completament $2x^3-8x$ i troba'n totes les arrels.\par
\end{minipage}\par\vspace{1.6mm}
\par\noindent\begin{minipage}{\linewidth}
\textbf{34.} \emph{(PAU 2022)} Factoritza el denominador i indica el domini de $f(x)=\dfrac{9}{x^2+x-2}$.\par
\end{minipage}\par\vspace{1.6mm}
\par\noindent\begin{minipage}{\linewidth}
\textbf{35.} \emph{(PAU 2023)} Resol per factorització l'equació $4x^2-10x-6=0$.\par
\end{minipage}\par\vspace{1.6mm}
\par\noindent\begin{minipage}{\linewidth}
\textbf{36.} \emph{(PAU 2024)} Factoritza completament i resol $x^3-7x^2+6x=0$.\par
\end{minipage}\par\vspace{1.6mm}
\par\noindent\begin{minipage}{\linewidth}
\textbf{37.} \emph{(PAU 2022, guiat)} Comprova, factoritzant, que $x^3-3a^2x+2a^3=(x-a)^2(x+2a)$.\par
\end{minipage}\par\vspace{1.6mm}
\par\noindent\begin{minipage}{\linewidth}
\textbf{38.} \emph{(PAU 2025)} Resol:\par

\begin{itemize}[leftmargin=0pt,itemsep=3pt,topsep=3pt,parsep=0pt,label={}]
\item \textbf{a)} $x^3+1=0$ (dona només la solució real)
\item \textbf{b)} si $a^3=8$, troba el valor real de $a$
\end{itemize}
\end{minipage}\par\vspace{1.6mm}
\par\noindent\begin{minipage}{\linewidth}
\textbf{39.} \emph{(PAU 2025)} Factoritza completament $4p^3-5p^4$ i troba'n totes les arrels.\par
\end{minipage}\par\vspace{1.6mm}
\par\noindent\begin{minipage}{\linewidth}
\textbf{40.} \emph{(PAU 2025)} Donada $f(x)=\dfrac{x^2-2x}{x-1}$:\par

\begin{itemize}[leftmargin=0pt,itemsep=3pt,topsep=3pt,parsep=0pt,label={}]
\item \textbf{a)} Factoritza el numerador i troba els punts de tall amb l'eix $X$.
\item \textbf{b)} Comprova si existeix cap valor de $x$ que compleixi $x^2-2x+2=(x-1)^2$.
\end{itemize}
\end{minipage}\par\vspace{1.6mm}
\par\noindent\begin{minipage}{\linewidth}
\textbf{41.} \emph{(PAU 2024)} Resol l'equació $3x^2+2x+2=x^2+2x+4$ (passa-ho tot a un costat i factoritza).\par
\end{minipage}\par\vspace{1.6mm}
\par\noindent\begin{minipage}{\linewidth}
\textbf{42.} \emph{(PAU 2024)} Resol per factorització $x^2-2x-15=0$.\par
\end{minipage}\par\vspace{1.6mm}
\par\noindent\begin{minipage}{\linewidth}
\textbf{43.} \emph{(PAU 2025)} Troba els valors del paràmetre que anul·len cada expressió, factoritzant-la:\par

\begin{itemize}[leftmargin=0pt,itemsep=3pt,topsep=3pt,parsep=0pt,label={}]
\item \textbf{a)} $p^2-1$
\item \textbf{b)} $\lambda^2+3\lambda-10$
\end{itemize}
\end{minipage}\par\vspace{1.6mm}
\par\noindent\begin{minipage}{\linewidth}
\textbf{44.} \emph{(PAU 2026)} Troba els valors de $m$ que anul·len $m^2-4m+3$, factoritzant-lo.\par
\end{minipage}\par\vspace{1.6mm}
\vspace{1mm}\hrule\vspace{3mm}

\section*{Solucions}
\small\begin{multicols}{2}

\textbf{Bloc 1.}\par
\par\noindent\begin{minipage}{\linewidth}
\textbf{1.} a) grau $4$, coef. principal $3$, terme independent $-7$ · b) grau $3$, coef. principal $-5$, terme independent $-1$ · c) grau $5$, coef. principal $1$, terme independent $0$ · d) grau $0$, coef. principal $8$, terme independent $8$.\par
\end{minipage}\par\vspace{1.6mm}
\par\noindent\begin{minipage}{\linewidth}
\textbf{2.} a) $p(1)=0$: arrel · b) $p(2)=3$: no és arrel · c) $q(3)=0$: arrel · d) $q(-1)=4$: no és arrel.\par
\end{minipage}\par\vspace{1.6mm}
\par\noindent\begin{minipage}{\linewidth}
\textbf{3.} $p+q=x^3+2x^2+2x+4$; $p-q=3x^3-4x^2+6x-6$.\par
\end{minipage}\par\vspace{1.6mm}
\par\noindent\begin{minipage}{\linewidth}
\textbf{4.} a) $2x^3+5x^2-7x+2$ · b) $x^2+6x+9$ · c) $x^2-4$.\par
\end{minipage}\par\vspace{1.6mm}
\par\noindent\begin{minipage}{\linewidth}
\textbf{5.} a) quocient $2x-3$, residu $6x-8$ · b) quocient $x^2-1$, residu $0$.\par
\end{minipage}\par\vspace{1.6mm}
\par\noindent\begin{minipage}{\linewidth}
\textbf{6.} a) quocient $x^2-4x+3$, residu $0$ · b) quocient $2x^2-3x+1$, residu $0$ · c) quocient $x^3+x^2-4x-4$, residu $0$.\par
\end{minipage}\par\vspace{1.6mm}
\textbf{Bloc 2.}\par
\par\noindent\begin{minipage}{\linewidth}
\textbf{7.} a) $x(x-2)(x-3)$ · b) $2x(x-2)(x+3)$ · c) $x^2(x-3)(x+3)$ · d) $3x(x-2)(x+2)$.\par
\end{minipage}\par\vspace{1.6mm}
\par\noindent\begin{minipage}{\linewidth}
\textbf{8.} a) $(x-3)(x-1)(x+2)$ · b) $(x-2)(x+1)(x+3)$.\par
\end{minipage}\par\vspace{1.6mm}
\par\noindent\begin{minipage}{\linewidth}
\textbf{9.} a) $(x-1)(x^2-2x+2)$, $\Delta=-4$ · b) $(x-1)(x^2+2x+3)$, $\Delta=-8$.\par
\end{minipage}\par\vspace{1.6mm}
\par\noindent\begin{minipage}{\linewidth}
\textbf{10.} a) $(x-2)(x-\sqrt2)(x+\sqrt2)$ · b) $(x-1)(x-\sqrt3)(x+\sqrt3)$.\par
\end{minipage}\par\vspace{1.6mm}
\par\noindent\begin{minipage}{\linewidth}
\textbf{11.} a) $2(x-2)(x-4)(x+4)$ · b) $3(x-2)(x+1)(x+2)$.\par
\end{minipage}\par\vspace{1.6mm}
\par\noindent\begin{minipage}{\linewidth}
\textbf{12.} a) $(x-1)(x-2)(x+3)$ · b) $(x-2)(x^2-2x+2)$, $\Delta=-4$ · c) $(x+1)(x^2-3)$ · d) $2(x-3)(x+2)\left(x-\tfrac12\right)$ · e) $3(x-2)(x+1)\left(x-\tfrac23\right)$.\par
\end{minipage}\par\vspace{1.6mm}
\par\noindent\begin{minipage}{\linewidth}
\textbf{13.} $(x-3)(x-2)(x-1)(x+1)(x+2)$.\par
\end{minipage}\par\vspace{1.6mm}
\par\noindent\begin{minipage}{\linewidth}
\textbf{14.} $(x-2)(x-1)(x+1)(x+3)(x^2+x+1)$, $\Delta=-3$.\par
\end{minipage}\par\vspace{1.6mm}
\par\noindent\begin{minipage}{\linewidth}
\textbf{15.} $(x^2-2x+2)(x^2+2x+2)=x^4+4$ (comprova-ho desenvolupant); els dos factors tenen $\Delta=-4<0$, són irreductibles.\par
\end{minipage}\par\vspace{1.6mm}
\par\noindent\begin{minipage}{\linewidth}
\textbf{16.} a) MCD$=x-3$, MCM$=(x-3)(x+3)(x+2)$ · b) MCD$=x+1$, MCM$=x(x-1)(x+1)^2$.\par
\end{minipage}\par\vspace{1.6mm}
\par\noindent\begin{minipage}{\linewidth}
\textbf{17.} a) $x^3-4x^2+x+6$ · b) $x^3+4x^2+4x$.\par
\end{minipage}\par\vspace{1.6mm}
\textbf{Bloc 3.}\par
\par\noindent\begin{minipage}{\linewidth}
\textbf{18.} a) $\dfrac{x-1}{x+1}$ · b) $\dfrac{x+2}{x-3}$ · c) $\dfrac{x}{x-3}$ · d) $x+1$.\par
\end{minipage}\par\vspace{1.6mm}
\par\noindent\begin{minipage}{\linewidth}
\textbf{19.} a) F · b) C · c) F · d) C.\par
\end{minipage}\par\vspace{1.6mm}
\par\noindent\begin{minipage}{\linewidth}
\textbf{20.} a) $\dfrac{3x-1}{x^2-1}$ · b) $\dfrac{2}{x^2-4}$ · c) $\dfrac{x^2+3x-9}{x(x-3)}$.\par
\end{minipage}\par\vspace{1.6mm}
\par\noindent\begin{minipage}{\linewidth}
\textbf{21.} a) $\dfrac{x+1}{x+3}$ · b) $x+2$.\par
\end{minipage}\par\vspace{1.6mm}
\par\noindent\begin{minipage}{\linewidth}
\textbf{22.} $\dfrac{2x-1}{x(x+1)}$.\par
\end{minipage}\par\vspace{1.6mm}
\par\noindent\begin{minipage}{\linewidth}
\textbf{23.} a) i b): comprova operant el costat esquerre fins a obtenir el dret.\par
\end{minipage}\par\vspace{1.6mm}
\textbf{Bloc 4.}\par
\par\noindent\begin{minipage}{\linewidth}
\textbf{24.} $x^2-4=(x-2)(x+2)$.\par
\end{minipage}\par\vspace{1.6mm}
\par\noindent\begin{minipage}{\linewidth}
\textbf{25.} $x+3$.\par
\end{minipage}\par\vspace{1.6mm}
\par\noindent\begin{minipage}{\linewidth}
\textbf{26.} $x+3$; en $x=5$: $8$.\par
\end{minipage}\par\vspace{1.6mm}
\par\noindent\begin{minipage}{\linewidth}
\textbf{27.} $\dfrac{x(x+2)}{x+1}$.\par
\end{minipage}\par\vspace{1.6mm}
\par\noindent\begin{minipage}{\linewidth}
\textbf{28.} $\dfrac{x(x+1)}{x+3}$; no definida en $x=3$ i $x=-3$.\par
\end{minipage}\par\vspace{1.6mm}
\par\noindent\begin{minipage}{\linewidth}
\textbf{29.} a) $V(x)=4x^3-64x^2+240x$ · b) $V(2)=256\,\text{cm}^3$.\par
\end{minipage}\par\vspace{1.6mm}
\par\noindent\begin{minipage}{\linewidth}
\textbf{30.} a) $B(x)=x^3-2x^2-5x+6$ · b) $B(x)=(x+2)(x-1)(x-3)$, punts d'equilibri $x=1$ i $x=3$ · c) $B(2)=-4$: pèrdues de $4\,000$ € en produir $200$ unitats.\par
\end{minipage}\par\vspace{1.6mm}
\par\noindent\begin{minipage}{\linewidth}
\textbf{31.} a) $N(5)=18$ (mil usuaris) · b) $N(t)=(t+1)(t-2)(t-4)$; amb $t\geq0$, el model val zero a $t=2$ i $t=4$ mesos.\par
\end{minipage}\par\vspace{1.6mm}
\textbf{Bloc 5.}\par
\par\noindent\begin{minipage}{\linewidth}
\textbf{32.} $x(x-2)^2$; arrels $x=0$ i $x=2$ (doble).\par
\end{minipage}\par\vspace{1.6mm}
\par\noindent\begin{minipage}{\linewidth}
\textbf{33.} $2x(x-2)(x+2)$; arrels $x=-2,0,2$.\par
\end{minipage}\par\vspace{1.6mm}
\par\noindent\begin{minipage}{\linewidth}
\textbf{34.} $x^2+x-2=(x-1)(x+2)$; domini $=\mathbb{R}\setminus\{-2,1\}$.\par
\end{minipage}\par\vspace{1.6mm}
\par\noindent\begin{minipage}{\linewidth}
\textbf{35.} $2(x-3)(2x+1)=0$; arrels $x=3$ i $x=-\tfrac12$.\par
\end{minipage}\par\vspace{1.6mm}
\par\noindent\begin{minipage}{\linewidth}
\textbf{36.} $x(x-1)(x-6)=0$; arrels $x=0,1,6$.\par
\end{minipage}\par\vspace{1.6mm}
\par\noindent\begin{minipage}{\linewidth}
\textbf{37.} $(x-a)^2(x+2a)=x^3-3a^2x+2a^3$ (comprova-ho desenvolupant).\par
\end{minipage}\par\vspace{1.6mm}
\par\noindent\begin{minipage}{\linewidth}
\textbf{38.} a) $x=-1$ · b) $a=2$.\par
\end{minipage}\par\vspace{1.6mm}
\par\noindent\begin{minipage}{\linewidth}
\textbf{39.} $p^3(4-5p)$; arrels $p=0$ (triple) i $p=\tfrac45$.\par
\end{minipage}\par\vspace{1.6mm}
\par\noindent\begin{minipage}{\linewidth}
\textbf{40.} a) $x(x-2)$; talls a $x=0$ i $x=2$ · b) l'equació es redueix a $2=1$: no té solució.\par
\end{minipage}\par\vspace{1.6mm}
\par\noindent\begin{minipage}{\linewidth}
\textbf{41.} $2x^2-2=0$; arrels $x=\pm1$.\par
\end{minipage}\par\vspace{1.6mm}
\par\noindent\begin{minipage}{\linewidth}
\textbf{42.} $(x-5)(x+3)=0$; arrels $x=5,-3$.\par
\end{minipage}\par\vspace{1.6mm}
\par\noindent\begin{minipage}{\linewidth}
\textbf{43.} a) $p=\pm1$ · b) $\lambda=-5,2$.\par
\end{minipage}\par\vspace{1.6mm}
\par\noindent\begin{minipage}{\linewidth}
\textbf{44.} $(m-1)(m-3)=0$; $m=1,3$.\par
\end{minipage}\par\vspace{1.6mm}
\end{multicols}
\end{document}
